\section{The original line trace inherits the proved structure-trace square}

\begin{lemma}[Line trace from an actual structure-trace square]
\label{lem:kn-native-line-trace-square-from-structure}
\lean{EnvelopeProperty.Vanishing.NadelGlobalCanonicalTrace.lineTrace_nativeSquare_of_structureSquare}
\uses{lem:kn-native-module-square-weighted-unit,lem:line-pf-unit-trace-native-local-coefficient,lem:scheme-native-coherent-tensor-ext}
\leanok
For any commuting scheme square and any original invertible base line,
suppose the actual structure-module morphisms $\tau,\nu$ satisfy their
structure-trace square. Then their line traces, formed using the SAME
original PFs and structure unitors, satisfy the native pulled-line
square. This is generic transport support; the Nadel consumer must
supply the previously proved original structure-trace factor internally.
\end{lemma}
\begin{proof}
Precompose with the original line PF followed by its original structure
left unitor; this is an isomorphism. Test every original pure tensor
$b\otimes m$ on every open. The left trace gives $\tau(b)m$, whose
native restriction is $s^\sharp(\tau(b))\eta_s(m)$. The original
weighted unit square sends the right input to
$J(b)\eta_q(\eta_s(m))$, and the original inverse-PF local reader
gives $\nu(J(b))\eta_s(m)$. The assumed structure-module equation
equates precisely these two scalar coefficients. Native tensor
extensionality and cancellation of the original isomorphism finish.
\end{proof}
